% Immutable cited-input supplement. Target proofs remain embargoed.
% Target: night:P4.37; tag 02XA.
% Frozen source SHA256 74497bfe327fae25c9cdcb17334b3bf42a9793a9df7c0a4404d97568f2f760f0.
% CompendiumRef is resolved against the cited S6_09 source edition, not silently identified with the current series statement.

% Exact cited input: series:II / night:P4.23
\begin{proposition}[P4.23]\hypertarget{tag-02WV}{}\label{night:P4.23}\label{stmt:P4.23}


In the standing torus-fibration setting of this section, let an irreducible singular support
$S$ have smooth normalization $T$. If its minimal model $T_0$
belongs to class VII and $b_2(T_0)>0$, then $e(S)>0$.
\end{proposition}

% Exact cited input: series:II / night:P4.26
\begin{proposition}[P4.26]\hypertarget{tag-02WY}{}\label{night:P4.26}\label{stmt:P4.26}


In Part~4.2 every irreducible reduced fibre support
has nonnegative topological Euler number. For a
singular irreducible support, rational resolution
of the normalization gives strictly positive
Euler number; so does a resolution whose minimal
model is of class VII with positive second
Betti number. These assertions permit singular
normalization.

Every normal-crossing support has the
nonnegative even Euler number of P4.6.
Consequently, for a fibration all of whose
reduced supports are irreducible or have
normal crossings,
\[
 B=\mathbb P^1,\qquad c_3(Y)\in\{0,2\},
 \qquad b_3(Y)=2-c_3(Y).
\]
This last conclusion does not assume a
classification of the equality cases for
irreducible nonnormal supports.
\end{proposition}

% Exact cited input: series:II / night:P4.36
\begin{proposition}[P4.36]\hypertarget{tag-02X9}{}\label{night:P4.36}\label{stmt:P4.36}


Every irreducible singular support in Part~4.2 whose
normalization has a Hopf minimal model has smooth minimal
Hopf normalization and Euler number zero. The assertion
does not assume initially that the normalization is smooth
or that its Euler number vanishes.
\end{proposition}

